% capitulos/capitulado-04-conclusao.tex
%
% Capítulo final do MODELO POR CAPÍTULOS AUTÔNOMOS. Diferentemente
% das "Considerações Finais" do modelo tradicional, este capítulo
% integra explicitamente as contribuições de cada capítulo de
% conteúdo em torno de uma síntese única, sem repetir a metodologia.
%
% Todo o conteúdo abaixo é fictício e serve apenas para demonstrar o
% padrão estrutural. Substitua pelo conteúdo real da sua pesquisa.

\textbf{Impacto científico:} [Sintetize aqui como os achados dos
capítulos de conteúdo (\autoref{cap:capitulo2}, \autoref{cap:capitulo3}
etc.) se articulam em torno de uma cadeia causal ou conceitual única,
e qual contribuição original isso representa para a área
\cite{exemplo:artigo:2024,exemplo:capitulo:2019}.]

\textbf{Impacto tecnológico:} O arcabouço integrado pelos capítulos
anteriores entrega: (i) [primeira entrega, com remissão ao capítulo
correspondente, por exemplo \autoref{cap:capitulo2}]; (ii) [segunda
entrega, com remissão ao \autoref{cap:capitulo3}]; (iii) [demais
entregas relevantes, replicando o padrão acima para cada capítulo de
conteúdo].

\textbf{Impacto social:} [Descreva aqui os beneficiários diretos e
indiretos dos resultados — gestores públicos, comunidades, setores
produtivos — e como os resultados subsidiam decisões concretas]
\cite{brasil:lei:2012}.

\textbf{Grau de inovação:} [Avalie o grau de inovação do trabalho
(baixo/médio/alto), justificando com base na combinação inédita de
abordagens empregada nos capítulos de conteúdo.]

\textbf{Limitações e trabalhos futuros:} [Descreva as limitações
reconhecidas do desenho adotado e aponte direções de pesquisa que
permanecem em aberto para trabalhos futuros.]
